% 第 52 章　氢原子：量子力学的样板间
\chapter{氢原子：量子力学的样板间}

第~51~章留下了一个悬案：卢瑟福发现原子是一个极小的核加上外面的电子，可按照经典电磁学，这样的电子应该在极短的时间内把能量辐射光，螺旋着栽进核里——原子根本不该稳定存在。玻尔用几条硬性规定强行稳住了原子，居然还算准了氢光谱的波长，但那些规定本身毫无理由。本章把悬案结案：用第~43~章的量子态和第~44~章的演化方程，从头解出氢原子。我们将看到，能级、轨道形状、跃迁规则不再是规定，而是同一个方程自然长出来的结果。氢原子之所以被称为量子力学的“样板间”，是因为它是自然界里唯一一个可以几乎完全精确求解、又和实验对得严丝合缝的真实系统——新理论的每一颗螺丝钉，都可以在这里拧一拧、验一验。

\step{反常识开场}

找一根装氢气的放电管（物理实验室的常见教具，样子像一根细长的霓虹灯管），两端加上高电压，氢气被电流击穿，发出一种柔和的粉红色光。霓虹灯发红橙色、氩气灯发蓝紫色、氢气灯发粉红色，这本身已经值得想一想：气体发光的颜色，为什么由气体的种类决定，而不是由电压高低决定？

更奇怪的在后面。把这束粉红色的光送进三棱镜，或者让它经过一片衍射光栅，你看到的不是一道连续铺开的彩虹，而是几条孤零零的亮线，像黑夜背景上的一张条形码：一条明亮的红线，一条青绿色的线，一条蓝紫色的线，再往里还有几条更暗的紫线，然后——没有了。其余所有颜色、所有波长，一概缺席。太阳光经过同样的光栅，铺开的倒是一道完整彩虹，但彩虹上嵌着几百条细细的暗线；其中几条暗线的位置，和氢放电管亮线的位置分毫不差。

一八八五年，瑞士中学数学教师巴尔末盯着氢的四条可见亮线的波长，发现了一个近乎不可思议的规律：把波长代入一个只含整数的简单公式，\(\frac{1}{4}-\frac{1}{n^2}\)（\(n\) 取 3、4、5、6），再乘上一个常数，四条线的位置一条不差地全部命中，误差小于千分之一。请注意这意味着什么：原子内部没有任何齿轮和刻度盘，可它发光的波长里藏着整数。原子在“数数”。它数的是什么？为什么全世界每一颗氢原子——实验室里的、太阳上的、几十亿光年外星系里的——数出来的是同一套数？

\step{先猜一猜}

把你的直觉预测写下来，越具体越好。

第一题。假如第~51~章的行星图像是真的：电子像一颗小行星，沿椭圆轨道绕核飞行，靠电磁力维系。它如果一边飞一边以光的形式漏掉能量，轨道就会越缩越小，转得越来越快。你猜它发出的光，颜色（也就是频率）会怎么变？是固定不变的几个颜色，还是一道平滑爬升的连续彩虹？

第二题。巴尔末公式里出现了整数。在你的经验里，什么东西天然和整数挂钩？一根两端固定的吉他弦只能弹出基频和整数倍泛音；一排台阶只有第几级、没有第几级半。你猜：氢光谱里的整数，是巧合，还是原子内部真的存在某种“一格一格”的结构？如果存在，格子是谁画出来的？

第三题。流水线上生产的零件再精密也有公差，两只同样型号的手表走时速率也有微小差别。可实验上，地球上的一颗氢原子和某颗遥远恒星上的一颗氢原子，发光波长一模一样，差别小到测不出来。你猜：原子为什么能做到“零公差”？它们之间是靠什么对齐的？

写好后夹在书页里。本章结束时我们逐条对答案。

\step{动手实验}

\begin{experiment}[用旧光盘给光“查户口”]
\textbf{材料}：一张废弃的 CD 或 DVD（它的刻痕就是一片现成的衍射光栅）；几种光源——白色手机屏幕、白光 LED 台灯、日光型荧光灯或节能灯、白炽灯泡或蜡烛，能透过窗户看到的霓虹招牌更好。注意：任何时候不要用光盘反射着直视太阳。

\textbf{第一步}：关掉屋里其他灯，点亮手机屏幕上一张纯白图片。把 CD 的无标签面斜对着屏幕，在某个角度上你会看到彩虹一样的光带——这是白光被光栅按波长摊开的颜色序列。慢慢调整角度，找到最亮、最清楚的一条光谱。

\textbf{第二步}：换一盏白炽灯或蜡烛，重复观察。你会看到一道平滑的彩虹，从红到紫连续铺开，没有任何空缺。

\textbf{第三步}：换成荧光灯管或节能灯。仔细看：在连续而暗淡的彩虹背景上，叠加着几条格外明亮的细线——一条亮绿线尤其醒目（它来自灯管里的汞），另外还能认出红、蓝的几条。这就是“条形码”叠在“彩虹”上的样子。

\textbf{第四步}：换成白光 LED，比较它和荧光灯的光谱：LED 通常是蓝光区域一个亮峰，加上黄绿区域一个宽宽的隆起，缺少荧光灯那种刀锋一样的细线。

\textbf{记录与思考}：把四种光源的“光谱形状”各画一笔：连续彩虹、彩虹加亮线、驼峰。请回答：同样都叫做“白光”，为什么光栅底下现出三种完全不同的原形？光源发光的颜色清单，到底是由什么决定的？
\end{experiment}

这个实验里你亲手验证了本章最关键的一个事实：气体原子发光是“挑波长”的，它们只发出一张离散的颜色清单。接下来要回答的是：这张清单是谁开的？

\step{旧模型遇到麻烦}

先把第~51~章的结局重演一遍，但这次把账算成实数。卢瑟福的原子里，电子绕核运动。经典电磁学有一条非常硬的结论：电荷做加速运动时必然向外辐射电磁波，带走能量。绕核转动的电子时刻都在拐弯，拐弯就是加速，所以它必须不停地发光。光发得越多，电子能量越低，轨道半径越小，转得越快——辐射反而更猛。这是一个加速失控的死亡螺旋。用经典公式估算，一颗初始尺寸约 \(10^{-10}\)~m 的氢原子，电子从出发到栽进核里，全过程只要大约 \(10^{-11}\)~s——一百亿分之一秒。而现实中的氢原子已经存在了一百三十八亿年，从宇宙诞生后不久一直活到今天。这是物理学史上理论预言与实验观测之间最悬殊的冲突之一：差了二十九个数量级，比“预言一只蚊子的寿命和银河系同龄”还要离谱。

更糟的是，就算电子坠得慢一点，死亡螺旋预言的光谱也不对。轨道平滑收缩，转动频率平滑升高，辐射的光就该是一道连续爬坡的彩虹——可实验给你的是几条孤零零的亮线。行星模型不但解释不了条形码，它预言的恰恰是条形码的反面。

玻尔在一九一三年的办法是直接立法：规定电子只能待在某些特殊轨道上，待在上面就禁止辐射；只有从一个允许轨道“跳”到另一个允许轨道时，才把能量差以一颗光子的形式吐出来或吞进去。这套规定出人意料地成功：由它推出的能量公式 \(E_n=-13.6/n^2\)（单位 eV）分毫不差地复现了巴尔末的整数规律，还预言了后来果然被找到的紫外线系和红外线系。但请诚实地清点玻尔模型的欠账：第一，“只允许某些轨道”“定态禁止辐射”都是硬贴上去的规定，没有任何更深的理由；第二，换到有两个电子的氦原子，这套规定就算不动了，更重的原子更是一筹莫展；第三，它说不清每条谱线为什么有明有暗——清单对得上，账目的粗细对不上；第四，精密光谱仪下每条氢线其实是一簇靠得极近的细线，玻尔模型对此完全失语。一个好补丁能把机器重新发动，但补丁不会变成发动机。我们需要一个让“格子”自己长出来的理论。

\step{发明新概念}

手头的工具其实已经备齐了。第~43~章告诉我们：一个电子的状态不是“某时刻在某处以某速度飞行”，而是一个量子态，它给出各种测量结果的概率幅；第~44~章告诉我们：量子态随时间的演化由薛定谔方程决定，方程里只需写明电子所处环境的势能。氢原子的环境再简单不过：电子在原子核的库仑势里运动，势能 \(V(r)=-\dfrac{e^2}{4\pi\varepsilon_0 r}\)，像一只向中心无限加深的圆碗。剩下的问题只剩一个：把量子态放进这只碗里，方程允许什么样的状态长期存在？

先用老朋友找直觉。两端固定的吉他弦上，不是所有的波都能留下来：波长必须恰好让弦上放下整数个“半波”，其余的波形刚出现就自相抵消。能留下来的驻波，频率天然是一份一份离散的——整数就是这样进入物理学的。原子里的电子与吉他弦有相似之处，也有必须说清楚的不同之处。相似之处在于：量子态的波被关在一个有限区域里（碗壁把它约束在核附近），同样只有某些“花样”能首尾相接地自洽存在，其余的在方程里根本立不住——能级离散，不是谁的规定，而是“波必须自我闭合”这个条件的直接后果。不同之处在于：吉他弦的波是弦这种物质在上下振动，而量子态的波是概率幅的分布，没有什么介质在动，也不能把电子想象成一小团被抹开的云雾物质；弦是一维的，原子是三维的，所以花样要丰富得多。

三维碗里的驻波花样，需要三个“编号”才能标定，就像鼓面的振动模式要用两个整数编号一样。第一个编号叫主量子数 \(n\)，取 1、2、3……它大致决定花样有多大、能量有多高：\(n\) 越大，概率分布伸得越远，能量越接近零（越自由）。第二个叫角量子数 \(l\)，取 0 到 \(n-1\)，它标记花样的“转动类别”：\(l=0\) 的花样是球对称的一团（记作 s），\(l=1\) 的花样是哑铃形的两瓣（记作 p），\(l=2\) 是更复杂的四瓣（记作 d）。第三个叫磁量子数 \(m\)，取 \(-l\) 到 \(+l\) 的整数，标记同一个转动类别在空间的取向：哑铃可以朝东、朝西、朝上，共 \(2l+1\) 种取向。再加上第~47~章的自旋——电子自带的一个没有经典对应物的内禀“角动量标签”，只有两种取值——一颗电子在氢原子里的“座位”就由四个数 \((n,l,m,m_s)\) 完全标定。

这里必须踩一脚刹车，纠正一个几乎人人带在脑子里的图像。“轨道”（orbital）是历史遗留下来的一个词，它指的不是一条轨迹，而是上面说的某一种空间花样——一个量子态。对于 \(n=1,\,l=0\) 的基态，量子态给出的不是“电子在某个球面上跑圈”，而是一个球对称的概率分布：问“位置测量可能在哪些地方找到电子”，答案是越靠近核的区域概率密度越高，但把概率乘上球壳的大小之后，“最容易找到电子的半径”出现在约 \(5.3\times10^{-11}\)~m 处。这颗电子此刻不“在”任何地方；不是它藏在哪里我们没本事知道，而是“确定位置”这个属性在测量之前对它根本不适用——这不是仪器粗糙的问题，而是量子态本身的性质（第~45~章已经仔细谈过这一点）。把轨道画成行星式的确定路线，等于把账簿当成了货物。

还剩最后一个概念障碍：电子既然在核附近“加速”，基态为什么不辐射、不坍缩？旧账本是这么记的：加速电荷必辐射。新账本改写了一条更基本的规则：处于定态的量子态，其概率分布不随时间变化——就像吉他弦的驻波，弦上每一点都在动，但波形的轮廓纹丝不动。一个概率分布完全静止的状态，不产生任何随时间变化的电流和电荷分布，按电磁学本身就没有辐射可言。坍缩之所以被顶住，靠的是第~45~章的老朋友：不确定关系。把电子往核上压，位置越收越窄，动量的涨落就被迫越大，动能代价猛涨。收缩要付的动能账和靠近核省下的势能账互相顶牛，平衡点就在 \(5.3\times10^{-11}\)~m 附近——这就是原子的尺寸。原子不坍缩，不是因为有根看不见的支柱，而是因为“再小下去更贵”。

\begin{oneline}
氢原子是一只三维的“共鸣箱”：电子的量子态在这只箱子里只允许一整套离散的驻波花样，每种花样用四个量子数编号、对应一个确定的能量；原子发光不是电子沿轨道滚下楼梯，而是量子态从一种花样换成另一种花样，能量差由一颗光子整份带走——光谱的条形码，就是这只共鸣箱的价目表。
\end{oneline}

\step{数学翻译器}

解库仑势中的薛定谔方程，定态能量干净得惊人：
\[
E_n=-\frac{13.6\ \mathrm{eV}}{n^2},\qquad n=1,2,3,\dots
\]
左边回答的问题：电子被束缚在核附近时，总能量可以取哪些值。右边只有两样东西：一个由基本常数（电子质量、电荷、普朗克常数、真空介电常数）算出来的能量单位 13.6~eV，和整数 \(n\) 的平方。负号的意思是“束缚”：能量为零对应电子恰好挣脱原子核，待在原子里的状态都要从外界“借”来能量才能解放，所以记为负。

按老规矩，立刻用特殊情况检查这条公式。\(n=1\)：基态能量 \(-13.6\)~eV，意思是把氢原子的电子彻底剥掉至少需要 13.6~eV——实验测出的氢电离能是 13.598~eV，对上了。\(n\to\infty\)：能量趋于零，电子越来越远、越来越自由，公式平滑地衔接上“自由电子能量非负”的经典世界，合理。\(n=0\) 或 \(n=-2\)：公式不许代入，量子数从 1 开始且只取正整数，这不是数学抠门，而是方程里根本没有这样的自洽花样。能量间隔：\(n=1\) 到 \(n=2\) 差 10.2~eV，\(n=2\) 到 \(n=3\) 只差 1.9~eV，越往上越挤——能级图下面宽、上面密，像一架踏板越到高处越密的梯子，这个形状直接决定了光谱的样子。

发光的规则只有一条：能量守恒。量子态从 \(n_i\) 换到 \(n_f\)（\(n_i>n_f\)），能量差变成一颗光子：
\[
E_\gamma=E_{n_i}-E_{n_f}=13.6\ \mathrm{eV}\times\left(\frac{1}{n_f^2}-\frac{1}{n_i^2}\right),
\qquad \lambda=\frac{hc}{E_\gamma}.
\]
代入 \(n_i=3,\,n_f=2\)：\(E_\gamma=13.6\times(\frac14-\frac19)=1.89\)~eV，用 \(hc\approx1240\ \mathrm{eV\cdot nm}\) 换算，\(\lambda\approx656\)~nm——正是氢放电管里那条明亮的红线（实测 656.3~nm）。代入 \(n_i=2,\,n_f=1\)：\(E_\gamma=10.2\)~eV，\(\lambda\approx122\)~nm，在紫外线区，肉眼看不见——这解释了为什么我们能用眼睛看到的氢线全是落到 \(n=2\) 的“巴尔末系”，而落到 \(n=1\) 的“莱曼系”全部藏在紫外。再检查两个极端：\(n_i=n_f\)，能量差为零，没有光子，等价于“没换花样不发光”，自洽；\(n_i\) 和 \(n_f\) 都很大且相邻，能级挤得极密，光子能量趋于零、频率趋于连续，量子世界平滑地融化回经典彩虹——这个“大量子数回到经典”的极限行为，是新理论必须通过的体检，它通过了。

\begin{mathbridge}
不用解方程，只用不确定关系就能估出基态的尺寸和能量。把电子大致限制在半径 \(r\) 的范围内，动量涨落至少是 \(\Delta p\sim \hbar/r\)，动能代价约为 \(\hbar^2/(2mr^2)\)；靠近核省下的库仑势能是 \(-e^2/(4\pi\varepsilon_0 r)\)。总账：
\[
E(r)\approx\frac{\hbar^2}{2mr^2}-\frac{e^2}{4\pi\varepsilon_0 r}.
\]
\(r\) 缩小时，第一项按 \(1/r^2\) 猛涨，第二项只按 \(1/r\) 变负——小到一定程度，收缩必然亏本。对 \(r\) 求最小值（令导数为零），得
\[
r_{\min}=\frac{4\pi\varepsilon_0\hbar^2}{me^2}\approx 5.3\times10^{-11}\ \mathrm{m},
\qquad E(r_{\min})\approx-13.6\ \mathrm{eV}.
\]
两个数字都命中精确解，这不是巧合：基态正是“收缩动能”与“聚拢势能”精打细算后的最小总账。这个估算同时说明了原子为什么是这个尺寸——普朗克常数就写在原子的半径里。
\end{mathbridge}

\begin{challenge}
三个量子数从哪里来？把球坐标下的薛定谔方程按“半径部分”和“角度部分”分离变量：要求角度部分的解绕一圈后严丝合缝地接回自己，就逼出两个整数——\(l\)（角动量大小满足 \(L^2=l(l+1)\hbar^2\)）和 \(m\)（角动量沿某方向的分量为 \(m\hbar\)）；再要求半径部分在无穷远处不发散，逼出 \(n\)，且要求 \(n\ge l+1\)。这就是“\(l\) 从 0 数到 \(n-1\)、\(m\) 从 \(-l\) 数到 \(+l\)”的由来：不是人为规定，而是“解必须自洽”的边界条件。由此还能数出每个 \(n\) 对应 \(n^2\) 种花样（不计自旋）。更妙的是谱线的明暗：哪种跃迁允许、哪种禁戒，由“选择定则”决定——单光子跃迁要求 \(\Delta l=\pm1\)。它的根源不是实验人员的任意禁令，而是对称性：光子带走一份角动量，初态和末态的转动花样必须恰好差一个角动量单位，角动量才守恒；否则振幅严格抵消，跃迁概率为零。对称性直接决定什么光能被看见——这是诺特定理精神在原子里的演出。
\end{challenge}

\begin{figure}[htbp]
\centering
\begin{tikzpicture}[>=Stealth,scale=1.0]
  % 能级：y = -13.6/n^2 eV，比例 0.25 cm/eV
  \draw[thick] (0,0) -- (6.4,0) node[right] {\(0\)（自由电子）};
  \draw[thick] (0,-0.136) -- (6.4,-0.136) node[right] {\(n{=}5\)};
  \draw[thick] (0,-0.213) -- (6.4,-0.213) node[right] {\(n{=}4\)};
  \draw[thick] (0,-0.378) -- (6.4,-0.378) node[right] {\(n{=}3\)};
  \draw[thick] (0,-0.85) -- (6.4,-0.85) node[right] {\(n{=}2\)};
  \draw[thick] (0,-3.4) -- (6.4,-3.4) node[right] {\(n{=}1\)（基态，\(-13.6\)~eV）};
  \draw[->,very thick,red!70!black] (1.6,-0.378) -- (1.6,-0.82);
  \node[right,red!70!black] at (1.7,-0.6) {H\(\alpha\)：656~nm 红光};
  \draw[->,very thick,blue!70!black] (3.4,-0.85) -- (3.4,-3.35);
  \node[right,blue!70!black] at (3.5,-2.1) {莱曼\(\alpha\)：122~nm 紫外};
  \draw[->,very thick,green!50!black,dashed] (5.3,-3.35) -- (5.3,-0.05);
  \node[right,green!50!black,align=left] at (5.4,-1.9) {吸收 13.6~eV\\ = 电离};
\end{tikzpicture}
\caption{氢原子的能级梯子：越往上踏板越密。向下的箭头是发光（能量差由光子带走），向上的虚线箭头是吸收入射光或碰撞能量。每条谱线的波长都由两级的高度差唯一决定。}
\end{figure}

\step{模型的边界}

把假设摆到台面上。本章解出的氢原子用了五个近似：电子是非相对论的（用薛定谔方程而不是狄拉克方程）；原子核被当成不动的点电荷（真实质子有大小、有质量、也在微微晃动）；忽略了电子自旋与轨道花样之间的磁耦合；忽略了光与原子相互作用里更微妙的量子电动力学效应；整个原子不受外场打扰。在这个范围内，模型精确得近乎蛮横——能级算到小数点后好几位都和实验一致。超出这个范围，偏差就一个接一个地现身，而且每一次现身都指向更深一层的理论：高速运动的修正让每条谱线劈成一簇细线，这是精细结构，由第~49~章的相对论性量子理论解释；\(2s\) 和 \(2p\) 之间一个约一千兆赫的微小错位（兰姆移位）超出了狄拉克理论，逼出了量子电动力学；电子自旋与核自旋的微小相互作用劈出的超精细结构，给了射电天文那条著名的 21~cm 谱线。样板间的意义正在于此：它不是终点，而是一台灵敏的探针，理论每深一层，氢原子都最先露馅。

外加电场和磁场会让谱线移动和分裂（斯塔克效应、塞曼效应），这不算模型失效——把外场写进势能再解方程即可，分裂的方式还反过来验证了磁量子数 \(m\) 的存在。真正的边界在多电子一侧：氦原子有两个电子，电子之间互相排斥，方程再也无法精确求解，本章的严格公式只剩定性意义。多电子世界靠近似和数值计算打天下，那是第~53~章的战场。

再清点三个典型误用。其一，把“轨道”画成行星轨道——前文已拆，不重复。其二，以为跃迁是一条看得见的“跳楼轨迹”：量子态从 \(n=3\) 换成 \(n=2\)，中间没有一段“电子逐渐下坠”的旅程，整个过程约 \(10^{-8}\)~s 内完成，能说的只有“初态是这个花样、末态是那个花样、一颗光子被带走”，追问“电子中途走到哪儿了”是拿经典剧本套量子舞台。其三，以为只有被“看见”才算跃迁——跃迁由原子和电磁场的相互作用决定，发不发光子、发什么颜色的光子，跟有没有人观测毫无关系；星光穿越几十亿年无人的太空，一路都是这个规则。

\step{回头解释开场现象}

现在拆开开场的那张条形码。氢放电管里，电流不断把原子撞醒（激发到各种高能级），原子的量子态随后纷纷换花样：\(3\to2\) 换出 656~nm 的红光——这就是放电管粉红色的主成分；\(4\to2\) 换出 486~nm 的青绿光，\(5\to2\) 换出 434~nm 的蓝紫光，全都落在 \(n=2\) 这一级上，这就是肉眼可见的巴尔末系。为什么只有这几条、没有中间色？因为能级是离散的：原子里根本不存在“能量差等于 600~nm”的两种花样，600~nm 的光子自然无从发出。条形码不是原子的怪癖，而是“能级梯子”加“能量守恒”的直接推论。

巴尔末的整数也有了出处：整数 \(n\) 是三维驻波花样的编号，波长公式里的 \(\frac{1}{n_f^2}-\frac{1}{n_i^2}\) 不过是一把能量差换算成波长的算盘。那位中学教师在一八八五年找到的经验公式，其实是薛定谔方程解的先声——他提前四十年摸到了答案的形状。

“零公差”之谜一并解决：氢原子的能级完全由 \(e\)、\(m\)、\(\hbar\)、\(\varepsilon_0\) 这几个普适常数决定，方程里没有“第几颗原子”这个变量。所有氢原子解同一道方程，自然得到同一组能级、同一张价目表。零件有公差是因为宏观物体由亿万个原子拼成，拼法总有出入；而一颗氢原子的“图纸”就是方程本身，没有拼装的余地。这也解释了太阳光里的暗线：太阳表面炽热的连续光谱穿过较冷的外层大气时，氢原子把波长恰好等于自己价目表的光子成批吞掉——暗线的位置和实验室里放电管亮线的位置分毫不差，因为它们本来就是同一组能级的同一种交易。几十亿光年外星系的氢发出同样的条形码，是全宇宙共用同一套物理定律的最直接证据。

这套理解还立刻变成了工具。太阳物理学家用只放行 656~nm 的 H\(\alpha\) 滤镜给太阳拍照，日面上的耀斑、暗条在这单一颜色里纤毫毕现；射电天文学家接收星际氢云发出的 21~cm 电波（超精细能级之间自旋“翻面”放出的光子），画出银河系的旋臂地图——单个原子的这种跃迁平均要等上千万年才发生一次，可银河系里的氢原子多达 \(10^{66}\) 量级，每时每刻都有天文数字的原子在“翻面”，汇成一架射电望远镜轻易可闻的合唱；天文学家还靠整张条形码的整体搬家测星体的运动：所有谱线等比例移向红端，说明光源在远离我们。一张价目表，同时是天体的身份证、温度计和测速仪。

\step{脑洞实验与挑战题}

\begin{thought}[如果普朗克常数大一万倍]
原子的半径里写着 \(\hbar\)：上面的数学桥给出 \(r_{\min}\propto\hbar^2\)。假如有一天宇宙把 \(\hbar\) 调大，原子会跟着发胖，\(\hbar\) 变为原来的一万倍多一点（约 \(1.4\times10^4\) 倍），氢原子就会长到厘米量级——一颗“氢弹珠”放在手心。请推理三个后果：第一，不确定关系 \(\Delta x\,\Delta p\gtrsim\hbar\) 在日常尺度上将变得无法忽略，你扔出去的球还会出现可分辨的位置—动量“模糊”吗？第二，原子的能级公式 \(E_n\propto 1/\hbar^2\)（把 13.6~eV 用常数展开看看），\(\hbar\) 变大意味着能级变浅还是变深？用可见光光子还能不能激发这样的原子？第三，回头想想：我们的世界之所以“泾渭分明”——位置确定、颜色确定、物体有清晰轮廓——在多大程度上仅仅因为 \(\hbar\) 是个极小的数？
\end{thought}

换个方向做一次尺度体操，顺便体会“轨道”图像错得有多远。把氢原子放大到一座操场那么大（半径 100~m），按真实比例，原子核只有约 \(2\times10^{-5}\)~m——一粒尘埃的大小，而电子连“一颗尘埃”都不是：基态的量子态是整个操场范围内一团球对称的概率分布。所以“尘埃核外面一颗小沙粒沿跑道兜圈”的画面，在这张比例图里没有立足之地；正确的画面更像整个操场上空弥漫着一层静止的、各处浓淡不同的“可能性雾气”——当然它也不像雾，雾由分子组成，而这团“雾气”只是对位置测量结果的预言规则。

\begin{puzzles}
\item 一批氢原子全被激发到 \(n=3\) 的态，随后各自自发跃迁回基态。它们发出的光里，最多能出现几种不同的波长？如果初态是 \(n=4\) 呢？（提示：把能级画成楼层，列出所有“下楼路线”；注意有的原子走直达、有的走中转，路线不同光子不同。）
\item 氢的 \(2p\to1s\) 跃迁大约 \(10^{-8}\)~s 就完成，而 \(2s\to1s\) 的单光子跃迁几乎被禁绝，一个卡在 \(2s\) 态的原子平均要拖上约十分之一秒才“脱困”——慢了上千万倍。能量差明明几乎一样，凭什么一个能发、一个不能发？（提示：回忆挑战盒里的选择定则。光子不只是能量包裹，它还带走一份角动量；\(2s\) 和 \(1s\) 的转动花样一模一样，角动量差为零，单光子“卸不下货”。）
\item 某天文学家测到一颗恒星的光谱，发现其中氢的全套谱线——相对间隔与氢的条形码完全吻合——但整体向红端搬了家，比如 H\(\alpha\) 出现在 660~nm 而不是 656~nm。他能下什么结论？如果另一颗恒星的谱线间隔与任何已知元素的价目表都对不上，又有哪些可能？（提示：整张价目表等比例平移，对应光源整体靠近或远离；而“间隔对不上”则不是平移能解释的——想想谱线还可能被什么加宽、劈裂或遮盖。）
\end{puzzles}

下一章，我们把这颗孤零零的电子放进人多的房间：几十个电子挤在同一个核周围，谁也不许和别人坐同一个座位——从这条“不相容”的规矩里，整张元素周期表将自己长出来。
